
@@M0@@ sur la catégorie des @@M1@@-Modules gradués est compatible au produit tensoriel; il en est encore de même du foncteur obtenu sur les catégories dérivées.


2.4. Pour tout entier @@M2@@, et tout @@M3@@-Module @@M4@@, posons :
@@M5@@
Le foncteur @@M6@@ ainsi défini est exact à gauche. Comme la catégorie des @@M7@@-Modules a assez d'injectifs, on peut définir un foncteur dérivé :
@@M8@@

Si @@M9@@ est un @@M10@@-Module localement libre de type fini, il existe un entier @@M11@@ tel que pour @@M12@@, @@M13@@ soit @@M14@@-acyclique. Il suffit pour le voir de montrer que, pour @@M15@@ assez grand, @@M16@@ est acyclique pour les foncteurs @@M17@@. Comme la tensorisation par @@M18@@ est un foncteur exact, cela résulte immédiatement de la suite spectrale des foncteurs composés et de 1.3.

2.5. On se propose de construire pour tout @@M19@@-Module @@M20@@ un homomorphisme fonctoriel en @@M21@@ :
@@M22@@
Pour cela on définit la restriction de @@M23@@ aux @@M24@@ (2.3.2). Or par construction on a, en posant @@M25@@ :
@@M26@@
Pour définir un homomorphisme~:
@@M27@@
